%-----------------------------QUESTÃO------------------------------------
%Cód.: TF3p23q01

\question [\questaopt] Considere as afirmativas a seguir:\\
I. A direção do vetor campo elétrico, em determinado ponto do 
espaço, coincide sempre com a direção da força que atua sobre uma
carga de prova colocada no mesmo ponto.\\
II. Cargas negativas, colocadas em um campo elétrico, tenderão a se
mover em sentido contrário ao do campo.\\
III. A intensidade do campo elétrico criado por uma carga pontual é, em
cada ponto, diretamente proporcional ao quadrado da carga que o
criou e inversamente proporcional à distância do ponto à carga.\\
IV. A intensidade do campo elétrico pode ser expressa em newton/
coulomb.
São verdadeiras:
				
%------------------------------ALTERNATIVAS-----------------------------

		\begin{EnvUplevel}
			\begin{choices}
			
			\choice somente I e II;  
			\choice somente III e IV;
			\correctchoice somente I, II e IV;
			\choice todas;
			\choice nenhuma.
				
			\end{choices}
		\end{EnvUplevel}

%--------------------------SOLUÇÃO---------------------------------------
		
		\begin{EnvUplevel}
		\begin{solution}[2cm]
		{\footnotesize
	I. Verdadeira\\
	A direção da força e do campo elétrico são iguais. O sentido é que
	pode ser diferente.\\
	II. Verdadeira\\
	Em cargas de prova negativas, a força elétrica e o campo elétrico
	possuem a mesma direção e sentidos opostos.\\
	III. Falsa\\
	$$
	E=k\times\frac{\left| Q  \right|}{\left( d \right)^{2}} \\
	$$
	IV. Verdadeira
	$$
	F = \left|q \right| \times E 
	$$ $$
	E = \frac{F}{\left|q \right|}
	$$
	No SI, a unidade de $ E $ pode ser \si[per-mode=symbol-or-fraction]{\newton\per\coulomb}.
		}
		
		\end{solution}
		\end{EnvUplevel}
